% Part: computability
% Chapter: recursive-functions
% Section: pr-functions-computable

\documentclass[../../../include/open-logic-section]{subfiles}

\begin{document}

\olfileid{cmp}{rec}{cmp}
\olsection{توابع بازگشتی اولیه محاسبه‌پذیرند}

فرض کنید تابع~$h$ با بازگشت اولیه تعریف شده باشد:
\begin{eqnarray*}
h(\vec x, 0)     & = & f(\vec x) \\
h(\vec x, y + 1) & = & g(\vec x, y, h(\vec x, y)),
\end{eqnarray*}
و فرض کنید توابع~$f$ و~$g$ محاسبه‌پذیرند. (از~$\vec x$ به‌عنوان
اختصار $x_0$، \dots، $x_{k-1}$ استفاده می‌کنیم.) آنگاه آشکارا
$h(\vec x, 0)$ محاسبه‌پذیر است، زیرا همان~$f(\vec x)$ است که بنا بر
فرض محاسبه‌پذیر است. سپس $h(\vec x, 1)$ نیز محاسبه‌پذیر است، زیرا
$1 = 0 + 1$ و بنابراین $h(\vec x, 1)$ برابر است با
\begin{align*}
 h(\vec x, 1) & = g(\vec x, 0, h(\vec x, 0)) =  g(\vec x, 0, f(\vec x)).
\intertext{می‌توانیم به همین شیوه ادامه دهیم و محاسبه کنیم:}
h(\vec x, 2) & = g(\vec x, 1, h(\vec x, 1)) = g(\vec x, 1, g(\vec x, 0, f(\vec x)))\\
h(\vec x, 3) & = g(\vec x, 2, h(\vec x, 2)) = g(\vec x, 2, g(\vec x, 1, g(\vec x, 0, f(\vec x))))\\
h(\vec x, 4) & = g(\vec x, 3, h(\vec x, 3)) = g(\vec x, 3, g(\vec x, 2, g(\vec x, 1, g(\vec x, 0, f(\vec x)))))\\
& \vdots
\end{align*}
پس برای محاسبهٔ $h(\vec x, y)$ در حالت کلی، به‌ترتیب
$h(\vec x, 0)$، $h(\vec x, 1)$، \dots را محاسبه می‌کنیم تا به
$h(\vec x, y)$ برسیم.

بنابراین اگر توابع~$f$ و~$g$ محاسبه‌پذیر باشند، تعریف بازگشتی اولیه
تابع محاسبه‌پذیر تازه‌ای به دست می‌دهد. ترکیب توابع نیز، اگر~$f$ و
$g_i$ها محاسبه‌پذیر باشند، تابعی محاسبه‌پذیر به دست می‌دهد.

چون توابع پایهٔ $\Zero$، $\Succ$ و~$\Proj{n}{i}$ محاسبه‌پذیرند و
ترکیب و بازگشت اولیه از توابع محاسبه‌پذیر توابع محاسبه‌پذیر می‌سازند،
هر تابع بازگشتی اولیه محاسبه‌پذیر است.

\end{document}
